%==============================================================================
\section{SARSA (State-Action-Reward-State-Action)}
\label{sec:sarsa}
%==============================================================================

\begin{dinhnghia}[title={SARSA}]
\textbf{SARSA} (State-Action-Reward-State-Action) là biến thể của Q-learning trong đó \textbf{chính sách mục tiêu trùng chính sách hành vi} (on-policy). Hai cặp trạng thái--hành động liên tiếp $(s,a)$ và $(s',a')$ cùng phần thưởng tức thời $r$ khi chuyển từ $s$ sang $s'$ quyết định giá trị cập nhật $Q$ --- đó là lý do tên thuật toán.
\end{dinhnghia}

\begin{ynghia}[title={Ý tưởng học thử--sai}]
SARSA là kỹ thuật học \textbf{on-policy} hiệu quả: agent thực hiện hành động, quan sát hậu quả, rồi điều chỉnh kế hoạch theo kết quả. Cốt lõi là \textbf{thử và sai}: chọn hành động trong tình huống, nhận phản hồi từ môi trường, cập nhật chiến lược dần theo chuỗi S--A--R--S--A.
\end{ynghia}

\subsection{Ví dụ robot trong mê cung}

\begin{vidu}[title={Robot đi mê cung}]
Giả sử dạy robot đi qua mê cung. Vị trí hiện tại là \textbf{trạng thái} $s$; mục tiêu tìm đường tốt tới đích. Mỗi bước robot có thể \textbf{hành động} (đi lên, xuống, trái, phải). Môi trường trả \textbf{phần thưởng} dương hoặc âm để báo hiệu bước đi tốt hay tệ. Sau khi thực hiện $a$ tại $s$, robot chuyển sang \textbf{trạng thái kế} $s'$; hành động \textbf{tiếp theo thực sự} $a'$ (theo chính sách hiện tại, ví dụ $\varepsilon$-greedy) sẽ xuất hiện trong công thức cập nhật SARSA --- khác với Q-learning dùng $\max_{a'} Q(s',a')$.
\end{vidu}

\tsimg{05_sarsa/img01.jpg}

\subsection{Phương trình cập nhật}

\begin{dinhnghia}[title={Công thức SARSA}]
Giá trị $Q(s,a)$ được cập nhật theo mẫu TD on-policy:
\[
Q(s,a) \leftarrow Q(s,a) + \alpha\,\bigl(r + \gamma\, Q(s',a') - Q(s,a)\bigr)
\]
trong đó $\alpha$ là tốc độ học, $\gamma$ hệ số chiết khấu, $r$ phần thưởng nhận được sau khi thực hiện $a$ tại $s$, $s'$ trạng thái kế, và $a'$ là \textbf{hành động thực sự} sẽ chọn tại $s'$ theo chính sách hành vi (không phải hành động tối ưu giả định).
\end{dinhnghia}

\begin{ynghia}[title={Đọc nhanh công thức}]
Mục tiêu TD là $r + \gamma Q(s',a')$: phần thưởng tức thời cộng giá trị chiết khấu của cặp trạng thái--hành động \textbf{kế tiếp theo chính sách đang theo}. Sai số $r + \gamma Q(s',a') - Q(s,a)$ nhân $\alpha$ rồi cộng vào $Q(s,a)$.
\end{ynghia}

\subsection{Các thành phần cốt lõi}

\begin{dinhnghia}[title={Trạng thái $S$}]
Trạng thái phản ánh môi trường, chứa mọi thông tin về tình huống hiện tại của agent (ví dụ ô trong mê cung).
\end{dinhnghia}

\begin{dinhnghia}[title={Hành động $A$}]
Hành động là quyết định của agent trong trạng thái hiện tại; thực hiện $a$ gây chuyển từ $s$ sang trạng thái khác.
\end{dinhnghia}

\begin{dinhnghia}[title={Phần thưởng $R$}]
Phần thưởng là tín hiệu phản hồi tức thời từ môi trường sau hành động tại $s$; hướng agent biết hành động nào được ưu tiên trong bối cảnh đó.
\end{dinhnghia}

\begin{dinhnghia}[title={Trạng thái kế $S'$}]
Sau khi thực hiện $a$ tại $s$, agent vào trạng thái mới $s'$ --- môi trường đã thay đổi theo hành động vừa chọn.
\end{dinhnghia}

\subsection{Cách SARSA hoạt động}

\begin{phuongphap}[title={Các bước thuật toán SARSA}]
\begin{enumerate}
  \item \textbf{Khởi tạo bảng $Q$}: gán $Q(s,a)$ giá trị tùy ý (hoặc 0); chọn trạng thái ban đầu $s$ và hành động ban đầu $a$ (thường bằng $\varepsilon$-greedy trên $Q$ hiện tại).
  \item \textbf{Khai thác vs khám phá}: khai thác dùng $Q$ đã biết để tối đa phần thưởng ngắn hạn; khám phá thử hành động ít chắc chắn để tìm chiến lược tốt hơn về lâu dài.
  \item \textbf{Thực thi và phản hồi}: thực hiện $a$, nhận $r$, chuyển sang $s'$.
  \item \textbf{Cập nhật $Q$}: chọn $a'$ tại $s'$ theo chính sách hành vi; cập nhật $Q(s,a)$ bằng công thức SARSA ở trên.
  \item \textbf{Lặp}: đặt $s \leftarrow s'$, $a \leftarrow a'$; lặp đến khi episode kết thúc. Qua nhiều vòng, $Q$ tiệm cận giá trị kỳ vọng theo \textbf{chính sách đang theo}, giúp quyết định ổn định hơn trong môi trường rủi ro.
\end{enumerate}
\end{phuongphap}

\begin{tinhchat}[title={Đặc điểm on-policy}]
Mỗi bước học phụ thuộc hành động \textbf{thực sự} sẽ làm ở $s'$, nên khám phá ảnh hưởng trực tiếp tới cập nhật --- phù hợp khi cần học an toàn, ổn định theo chính sách triển khai (y tế, giao thông, học cá nhân hóa).
\end{tinhchat}

\subsection{SARSA so với Q-learning}

\begin{vidu}[title={Bảng so sánh}]
\begin{tabularx}{\linewidth}{>{\raggedright\arraybackslash}p{3.4cm}XX}
\textbf{Tiêu chí} & \textbf{SARSA} & \textbf{Q-learning} \\
\midrule
Loại chính sách &
On-policy (mục tiêu = hành vi). &
Off-policy (học tối ưu trong khi hành vi có thể khác). \\
Quy tắc cập nhật &
$Q(s,a) \leftarrow Q(s,a) + \alpha\bigl(r + \gamma Q(s',a') - Q(s,a)\bigr)$ &
$Q(s,a) \leftarrow Q(s,a) + \alpha\bigl(r + \gamma \max_{a'} Q(s',a') - Q(s,a)\bigr)$ \\
Hội tụ &
Thường chậm hơn tới chính sách tối ưu toàn cục nhưng bám sát chính sách an toàn. &
Thường nhanh hơn tới $Q^*$ nếu đủ khám phá. \\
Khám phá--khai thác &
Khám phá ảnh hưởng trực tiếp bước học. &
Có thể tách chính sách khám phá khỏi chính sách học. \\
Cập nhật policy &
Dựa trên hành động \textbf{đã thực sự} chọn ở $s'$. &
Giả định hành động \textbf{tốt nhất} ở $s'$ ($\max_{a'}$). \\
Ứng dụng &
Môi trường cần ổn định, rủi ro cao khi lệch chính sách. &
Game, robot, giao dịch khi ưu tiên hiệu quả tối đa. \\
\end{tabularx}
\end{vidu}

\begin{luuy}[title={Chọn thuật toán}]
Dùng SARSA khi muốn học giá trị của \textbf{chính sách đang chạy} (kể cả $\varepsilon$-greedy). Dùng Q-learning khi muốn ước lượng $Q^*$ của chính sách tối ưu trong khi vẫn khám phá bằng chính sách khác.
\end{luuy}
